\documentclass[11pt]{article}
\usepackage[a4paper,margin=0.92in]{geometry}
\usepackage[T1]{fontenc}
\usepackage{lmodern}
\usepackage{amsmath,amssymb,amsthm,mathtools}
\usepackage{graphicx}
\usepackage{microtype}
\usepackage[hidelinks]{hyperref}
\usepackage{xcolor}
\usepackage{fancyhdr}

\definecolor{accent}{RGB}{47,79,106}
\newtheorem{theorem}{Theorem}
\newtheorem{lemma}[theorem]{Lemma}
\newcommand{\F}{\mathbb F}
\newcommand{\Inf}{\operatorname{Inf}}
\newcommand{\supp}{\operatorname{supp}}
\newcommand{\tr}{\operatorname{tr}}
\pagestyle{fancy}
\fancyhf{}
\setlength{\headheight}{14pt}
\lhead{Quantum $L^2$-KKL for commuting Pauli support}
\rhead{Lane 4}
\cfoot{\thepage}
\fancypagestyle{plain}{%
  \fancyhf{}%
  \lhead{Quantum $L^2$-KKL for commuting Pauli support}%
  \rhead{Lane 4}%
  \cfoot{\thepage}%
}

\title{\color{accent}Quantum $L^2$-KKL for Observables\\
with Commuting Pauli Support\\[3pt]
\large A complementary substantial partial result}
\author{Automated Mathematics Run \texttt{fa\_banach\_001}, lane 4}
\date{27 August 2026}

\begin{document}
\maketitle

\begin{center}
\fbox{\parbox{0.91\linewidth}{\textbf{Status.}
\texttt{candidate\_partial\_likely\_valid}.  The full quantum $L^2$-KKL
conjecture remains open.  This packet proves it for every balanced quantum
Boolean observable whose Pauli-Fourier support is contained in an abelian
Pauli subgroup.}}
\end{center}

\section{Source problem}

Blecher, Gao, and Xu, \emph{Geometric influences on quantum Boolean cubes},
arXiv:2409.00224, Section~7.1 on PDF page~35, state that the original
Montanaro--Osborne quantum KKL conjecture for $L^2$-influence remains open.

\begin{figure}[ht]
  \centering
  \includegraphics[width=\linewidth]{figures/open_problem_crop.png}
  \caption{The open-problem passage in Section 7.1 of the source paper.}
\end{figure}

In its balanced form, the conjecture asks whether every self-adjoint unitary
$A\in M_2(\mathbb C)^{\otimes n}$ with normalized trace $\tau(A)=0$ satisfies
\[
  \max_{1\leq j\leq n}\Inf_j^2(A)
  \geq c\frac{\log n}{n}
\]
for an absolute $c>0$.  Here $d_j$ removes the Pauli terms that are the
identity on qubit $j$, and $\Inf_j^2(A)=\lVert d_jA\rVert_2^2$.

\section{Statement and proof intuition}

Write $\mathcal P_n=\{I,X,Y,Z\}^{\otimes n}$ for the Hermitian Pauli basis,
with irrelevant signs handled separately.  A Pauli-supported observable is
called \emph{commuting} here if all Pauli strings with nonzero Fourier
coefficient lie in a common abelian Pauli subgroup.

The apparent obstacle is that a global Clifford diagonalizing that subgroup
need not preserve physical qubits, so classical KKL cannot simply be quoted in
the original coordinates.  The repair is linear-algebraic.  A rank-$m$
commuting Pauli subgroup is an $m$-dimensional isotropic subspace of
$\F_2^{2n}$.  Its $2n$ scalar $X/Z$ coordinate projections span the dual
space, so $m$ of them form a dual basis.  Relabel the Fourier characters by
that basis.  Classical KKL then finds a large influence in one chosen scalar
projection, and that event is contained in ``the Pauli is nonidentity on the
associated physical qubit.''  Thus the classical influence is bounded above
by a genuine quantum influence.

\begin{theorem}[Commuting-Pauli quantum KKL]\label{thm:main}
There is an absolute constant $c>0$ with the following property.  Let
$A=A^*$, $A^2=I$, and $\tau(A)=0$ on $n$ qubits.  If the Pauli-Fourier support
of $A$ is contained in an abelian Pauli subgroup, then
\[
  \max_{1\leq j\leq n}\Inf_j^2(A)
  \geq c\frac{\log(n+1)}{n}.
\]
Consequently the Montanaro--Osborne conjecture holds for this entire class.
\end{theorem}

\section{Proof}

Choose independent commuting Hermitian Pauli strings
$P_1,\ldots,P_m$ generating the subgroup spanned by the Fourier support, and
put
\[
 P_\xi=P_1^{\xi_1}\cdots P_m^{\xi_m},\qquad \xi\in\F_2^m.
\]
The generators may be chosen independent modulo scalar phases, so the
$P_\xi$ are distinct Pauli-basis elements and $m\leq n$.  The latter bound is
also the standard fact that an isotropic subspace of the symplectic space
$\F_2^{2n}$ has dimension at most $n$.

Expand
\[
 A=\sum_{\xi\in\F_2^m}a_\xi P_\xi.
\]
The commuting generators have all $2^m$ joint sign patterns, each with equal
multiplicity $2^{n-m}$.  Indeed, the joint spectral projector for
$x\in\{\pm1\}^m$ is
\[
 E_x=2^{-m}\prod_{r=1}^m(I+x_rP_r),
\]
and $\tr(E_x)=2^{n-m}$ because every nonidentity Pauli string is traceless.
Hence $A$ is the functional calculus of
\[
 f(x)=\sum_{\xi\in\F_2^m}a_\xi x^\xi.
\]
The conditions $A^2=I$ and $\tau(A)=0$ say respectively that
$f:\{\pm1\}^m\to\{\pm1\}$ is Boolean and balanced.  Moreover $a_\xi$ are
exactly its Walsh--Fourier coefficients.

Represent an element $P_\xi$ modulo phase by its binary symplectic vector
\[
 \rho(\xi)=(x_1(\xi),\ldots,x_n(\xi),z_1(\xi),\ldots,z_n(\xi))
 \in\F_2^{2n}.
\]
The linear map $\rho:\F_2^m\to\F_2^{2n}$ is injective.  Therefore its
$2n$ scalar coordinate functionals span $(\F_2^m)^*$.  Select a basis
$\ell_1,\ldots,\ell_m$ from these functionals.  For each $r$, let $j(r)$ be
the physical qubit from whose $X$- or $Z$-coordinate $\ell_r$ was selected.

The map
\[
 T:\F_2^m\longrightarrow\F_2^m,
 \qquad T\xi=(\ell_1(\xi),\ldots,\ell_m(\xi))
\]
is invertible.  An invertible linear change of Boolean variables induces the
dual relabeling of Walsh characters.  Thus there is a balanced Boolean
function $g$ on $m$ bits such that
\[
 \widehat g(T\xi)=a_\xi\qquad(\xi\in\F_2^m).
\]
For its $r$-th classical influence,
\begin{align*}
 \Inf_r(g)
 &=\sum_{\eta:\eta_r=1}|\widehat g(\eta)|^2
   =\sum_{\xi:\ell_r(\xi)=1}|a_\xi|^2\\
 &\leq \sum_{\xi:\,(x_{j(r)}(\xi),z_{j(r)}(\xi))\ne(0,0)}|a_\xi|^2
   =\Inf_{j(r)}^2(A).
\end{align*}
The inequality holds because a nonzero selected scalar coordinate forces the
Pauli string to be nonidentity on that qubit; the last equality is Pauli
orthogonality.

The classical KKL theorem, in a form valid for all $m\geq1$, gives an absolute
$c_0>0$ such that
\[
 \max_{1\leq r\leq m}\Inf_r(g)
 \geq c_0\frac{\log(m+1)}{m}\operatorname{Var}(g)
 =c_0\frac{\log(m+1)}{m}.
\]
Finally, $t\mapsto\log(t+1)/t$ decreases on $(0,\infty)$ and $m\leq n$.
Combining the last two displays proves Theorem~\ref{thm:main} with $c=c_0$.
\qed

\section{Verification and limitations}

The proof was checked in two independent ways.  First, the only dimension
count is the injective-map identity
$\operatorname{rank}(\rho)=\operatorname{rank}(\rho^*)=m$, which guarantees
the selected dual basis.  Second, the accompanying script
\texttt{code/verify\_commuting\_transfer.py} generates small random isotropic
Pauli subspaces and balanced Boolean functions and verifies the exact transfer
inequality $\Inf_r(g)\leq\Inf_{j(r)}^2(A)$ for every selected coordinate.  The
script is a consistency check, not part of the proof.

This result is complementary to the existing run packet for arXiv:2209.07279
on anticommuting Pauli support, which obtains the stronger order $n^{-1/2}$.
Neither
theorem handles a general unitary Pauli expansion, where commuting products
from different parts of the support can cancel in the equations imposed by
$A^2=I$.

\section{Novelty check and review recommendation}

A bounded search through 27 August 2026 used the source arXiv id and the exact
phrases ``commuting Pauli quantum KKL'', ``abelian Pauli influence KKL'',
``Pauli support KKL'', and ``stabilizer quantum KKL''.  It recovered the known
$L^p$, $p<2$, quantum KKL papers and this run's anticommuting partial result,
but no theorem matching the commuting-support reduction above.  Novelty
confidence is moderate rather than definitive: the proof is short enough that
it may be folklore through stabilizer diagonalization.

\paragraph{Human-review recommendation.}
Review as a likely-valid substantial partial result.  The highest-value checks
are: (i) that independent commuting Pauli generators realize all joint sign
patterns; (ii) that the $2n$ local scalar projections span the dual of their
isotropic subspace; and (iii) the direction of the influence comparison after
the Fourier-coordinate change.

\begin{thebibliography}{9}
\bibitem{BGX}
D. P. Blecher, L. Gao, and B. Xu,
\emph{Geometric influences on quantum Boolean cubes},
arXiv:2409.00224.

\bibitem{MO}
A. Montanaro and T. J. Osborne,
\emph{Quantum Boolean functions},
Chicago J. Theoret. Comput. Sci. (2010), arXiv:0810.2435.

\bibitem{KKL}
J. Kahn, G. Kalai, and N. Linial,
\emph{The influence of variables on Boolean functions},
Proc. 29th IEEE Symposium on Foundations of Computer Science (1988), 68--80.

\bibitem{RWZ}
C. Rouz\'e, M. Wirth, and H. Zhang,
\emph{Quantum Talagrand, KKL and Friedgut's theorems and the learnability of
quantum Boolean functions}, Comm. Math. Phys. 405 (2024), Article 95;
arXiv:2209.07279.
\end{thebibliography}

\end{document}
